%%%%%%
% Metadata
%%%%%%
\appendix
\chapter{Supplementary Materials and Implementation Details}

\renewcommand{\theequation}{A.\arabic{equation}}
\setcounter{equation}{0}

This appendix summarizes the key mathematical expressions and
implementation details used in the Alpha Recommendation System.
The formulas listed here provide a compact reference for the
computational procedures described in Chapter 4.

\section{Baseline Alpha Prediction}

The baseline significance level is predicted using a Random Forest
regression model trained on replication data.

\begin{equation}
	\alpha_{\text{base}} = f(N,d)
\end{equation}

where $N$ represents the sample size and $d$ denotes the expected
standardized effect size.

\section{Alpha Adjustment Pipeline}

The baseline alpha predicted by the model is adjusted through a
sequence of contextual transformations.

\subsection{Risk Tolerance Adjustment}

\begin{equation}
	\alpha' = \alpha_{\text{base}} \cdot k_1
\end{equation}

where $k_1$ represents the risk tolerance scaling factor.

\subsection{Test Direction Adjustment}

\begin{equation}
	\alpha'' =
	\begin{cases}
		2\alpha' & \text{if one-tailed}\\
		\alpha' & \text{if two-tailed}
	\end{cases}
\end{equation}

This adjustment accounts for whether the statistical test is
one-tailed or two-tailed.

\subsection{Multiple Testing Correction}

\begin{equation}
	\alpha_{\text{final}} = \frac{\alpha''}{m}
\end{equation}

where $m$ denotes the number of simultaneous hypothesis tests.

\section{Optimal Alpha Definition}

For studies in which the replication does not confirm the original
finding, a study-specific optimal significance level is derived using
the original reported p-value.

\begin{equation}
	\alpha_{\text{opt}} =
	\max \{\alpha \in A : \alpha < p_{\text{orig}}\}
\end{equation}

This rule selects the largest candidate alpha value that remains
strictly smaller than the reported p-value.

\section{Candidate Alpha Set}

The candidate set of significance levels used for selecting the
optimal alpha is defined as

\begin{equation}
	A =
	\{0.001,0.005,0.010,0.020,0.030,0.040,0.050,0.060,0.090,0.100\}
\end{equation}

Using a discrete set of candidate values ensures that the resulting
thresholds remain interpretable and comparable across studies.

\section{Statistical Recovery Formulas}

When original studies did not explicitly report sample size or effect
size, these quantities were reconstructed using standard statistical
relationships derived from reported test statistics.

\subsection{Effect Size from t-test}

\begin{equation}
	r = \sqrt{\frac{t^2}{t^2 + df}}
\end{equation}

\subsection{Sample Size from Degrees of Freedom}

\begin{equation}
	N = df + 2
\end{equation}

\subsection{Effect Size from Chi-Squared Test}

\begin{equation}
	\phi = \sqrt{\frac{\chi^2}{N}}
\end{equation}

\subsection{Effect Size from F-test}

\begin{equation}
	\eta^2 = \frac{F \cdot df_1}{F \cdot df_1 + df_2}
\end{equation}

\subsection{Conversion from $\eta^2$ to Correlation}

\begin{equation}
	r = \sqrt{\eta^2}
\end{equation}

\section{Model Evaluation Metric}

The performance of the Random Forest regression model is evaluated
using the Mean Absolute Error (MAE).

\begin{equation}
	MAE =
	\frac{1}{n}
	\sum_{i=1}^{n}
	|y_i - \hat{y}_i|
\end{equation}

where $y_i$ represents the observed optimal alpha values and
$\hat{y}_i$ denotes the predicted values.

\section{Python Implementation}

The Alpha Recommendation System was implemented in Python using the
\texttt{scikit-learn} machine learning library.

\begin{lstlisting}[language=Python]
	import pandas as pd
	from sklearn.ensemble import RandomForestRegressor
	
	data = pd.read_excel("master_alpha_dataset_updated.xlsx")
	
	X = data[["sample_size","effect_size"]]
	y = data["best_alpha"]
	
	model = RandomForestRegressor(n_estimators=200, random_state=42)
	model.fit(X, y)
\end{lstlisting}

\section{Source Code Repository}
\label{sec:gitappendix}
The full implementation of the Alpha Recommendation System,
including the trained model, dataset, and Streamlit application,
is available in the public GitHub repository:

\url{https://github.com/Abishek700/AlphaPredictionSystem}

The repository contains the dataset used for training, the Python
scripts for model training and evaluation, the serialized machine
learning model, and the Streamlit application used to generate
alpha recommendations.

\subsection{Running the Application Locally}

\begin{lstlisting}[language=bash]
	# 1. Clone the repository
	git clone https://github.com/Abishek700/AlphaPredictionSystem
	cd AlphaPredictionSystem
	
	# 2. Install dependencies
	pip install -r requirements.txt
	
	# 3. Launch the application
	streamlit run app.py
\end{lstlisting}

Once launched, Streamlit opens the interface automatically in the
default web browser, allowing users to enter study parameters and
generate alpha recommendations.